\section{Results}

% 纪律：本节所有数值均来自 validation/results/ 的真实输出（R 与 C++ 实跑）。

\subsection{Typical analysis workflow}

{\sloppy We applied the application to the bundled example dataset
(\texttt{examples/cq.csv}, \texttt{examples/design.csv}): the target gene
\textit{fos-glo-myc} against the reference gene \textit{Beta Actin} across
three treatment groups (0, 0.5, 1) with three biological replicates each.
With the $\Delta\Delta C_t$ method and Welch tests against the control,
relative expression rose to 1.92 at dose~0.5 and 2.74 at dose~1, both
significantly different from the control ($p = 0.006$ and $p = 0.022$); as
defined by Livak and Schmittgen the calibrator is the control group, which is
therefore $\approx 1$ (here $1.01 \pm 0.21$). These means, standard
deviations and $p$-values are reproduced without the application or the R
package by \texttt{validation/check\_example\_stats.R}, and the complete
output is shown in Supplementary Figure~S1B. On a four-fold dilution series
of four genes, the standard-curve module recovered slopes of $-2.07$ to
$-2.24$ $C_q$ per dilution step with $R^2 > 0.99$ for three of the four
primers (efficiencies 0.85--0.95; Supplementary Figure~S1A), illustrating the
efficiency estimation that, per MIQE guidance, determines which
quantification method is appropriate. The quality-control panel flagged the
single $C_q = 40$ well of the weakest primer, and the complete
workflow---import, configuration, analysis, export---is performed in the
graphical interface without writing code.\par}

\subsection{Equivalence to the R package}

{\sloppy We ran the C++ implementation and the original R package
\cite{li2022qpcrtools} on identical inputs for every shared
method---$\Delta\Delta C_t$ (\texttt{CalExp2ddCt}), the standard curve
(\texttt{CalCurve}) and the efficiency-based RqPCR method
(\texttt{CalExpRqPCR})---reading the same on-disk files, so any difference
reflects the implementations alone (Supplementary Table~S1 and Figure~S2).
For both expression methods, expression values, standard deviations and
standard errors matched to within floating-point precision (maximum relative
difference $< 4 \times 10^{-15}$). For the standard curve, slope, intercept,
$R^2$ and efficiency agreed to within $2.5 \times 10^{-3}$ relative
difference, the residual being attributable to the R package rounding its
tabulated output to two to four decimals while the underlying $C_q$ inputs
are identical; every point falls on the line $y = x$.\par}

The validation is fully reproducible from the open scripts in
\texttt{validation/} (see Availability). Preparing the port also revealed
defects in result aggregation that the R implementation shared---the input of
the reference-gene stability measure, replicate handling in the efficiency
method and the $\Delta\Delta C_t$ calibrator---as well as one specific to the
C++ code (multiple-comparison $p$-values); each was corrected and is covered
by a regression test.
